\section{Diseño del esquema y fragmentación}
\label{sec:diseno-esquema}

\subsection{Esquema físico}

La ubicación de la capa de datos dentro de la arquitectura completa del sistema —incluidos el
\texttt{api-gateway}, las aplicaciones cliente y \texttt{telemetry-service}, ya no solo los cinco
microservicios de la Entrega 2— se muestra en la Figura~\ref{fig:c4-nivel2} de la
Sección~\ref{sec:arquitectura-e4}. Esta sección se concentra en el diseño interno de esa capa: la
fragmentación del clúster de CockroachDB (rediseñada en esta entrega) y el pipeline de Spark que
la consume.

El diagrama entidad-relación completo de \texttt{ticket\_db} —la base propia de
\texttt{ticket-service} dentro del clúster de tres nodos de CockroachDB— se documenta en
\texttt{docs/diagrams/db-schema.md}. La tabla \texttt{tickets} es la de mayor cardinalidad y la
única fragmentada; \texttt{technicians} es una dimensión pequeña que no requiere fragmentación. La
Tabla~\ref{tab:esquema-tickets} detalla las columnas relevantes para la política de fragmentación.

\begin{table}[H]
\centering
\caption{Columnas relevantes de \texttt{tickets} tras el rediseño de la Entrega 3}
\label{tab:esquema-tickets}
\begin{tabular}{@{}llp{6.5cm}@{}}
\toprule
\textbf{Columna} & \textbf{Tipo} & \textbf{Rol} \\
\midrule
\texttt{created\_at} & \texttt{TIMESTAMPTZ} & Primer componente de la PK; columna de fragmentación \\
\texttt{id}           & \texttt{UUID}        & Segundo componente de la PK; respaldado además por índice único \texttt{tickets\_id\_key} \\
\texttt{zone}         & \texttt{STRING}      & Columna de negocio (RBAC de técnico, reportes); indexada, ya no es clave de fragmentación \\
\texttt{client\_id}   & \texttt{UUID}        & Identidad del cliente, resuelta a nivel de aplicación vía JWT de \texttt{auth-service} \\
\texttt{status}       & \texttt{STRING}      & Indexada (\texttt{idx\_tickets\_status}) \\
\bottomrule
\end{tabular}
\end{table}

\subsection{Política de fragmentación}

El fragmento del esquema en el Listado~\ref{lst:init-db} muestra la definición efectiva de la
tabla \texttt{tickets} tal como se aplica en la migración Flyway
\texttt{services/svc-principal/.../db/migration/V1\_\_init\_ticket\_schema.sql} (Entregable 5
de la guía de cierre; antes era un guion suelto sin versionar). La verificación de que las
migraciones se aplican de verdad no se limita a \texttt{ticket\_db}: además de
\texttt{TicketRepositoryIntegrationTest} (svc-principal), \texttt{auth-service} y
\texttt{report-service} tienen su propio \texttt{FlywayMigrationIntegrationTest} —cada uno
levanta un CockroachDB real vía Testcontainers, confirma que \texttt{flyway\_schema\_history}
registra la migración V1 como exitosa, y consulta su tabla principal—, en vez de depender solo
del arranque del \emph{stack} completo en el trabajo \texttt{integration} del flujo de CI.

\begin{lstlisting}[language=SQL, caption={Fragmentación por rango de \texttt{tickets} (extracto de \texttt{V1\_\_init\_ticket\_schema.sql})}, label={lst:init-db}]
CREATE TABLE IF NOT EXISTS tickets (
    created_at      TIMESTAMPTZ NOT NULL DEFAULT now(),
    id              UUID NOT NULL DEFAULT gen_random_uuid(),
    zone            STRING NOT NULL,
    client_id       UUID NOT NULL,
    technician_id   UUID REFERENCES technicians(id),
    category        STRING,
    priority        STRING,
    status          STRING NOT NULL DEFAULT 'NUEVO',
    description     STRING,
    sla_deadline    TIMESTAMPTZ,
    resolved_at     TIMESTAMPTZ,
    sla_breached    BOOL DEFAULT FALSE,
    PRIMARY KEY (created_at, id)
) PARTITION BY RANGE (created_at) (
    PARTITION tickets_2026_q1 VALUES FROM (MINVALUE) TO ('2026-04-01T00:00:00Z'),
    PARTITION tickets_2026_q2 VALUES FROM ('2026-04-01T00:00:00Z') TO ('2026-07-01T00:00:00Z'),
    PARTITION tickets_2026_q3 VALUES FROM ('2026-07-01T00:00:00Z') TO ('2026-10-01T00:00:00Z'),
    PARTITION tickets_2026_q4 VALUES FROM ('2026-10-01T00:00:00Z') TO (MAXVALUE)
);

CREATE UNIQUE INDEX IF NOT EXISTS tickets_id_key ON tickets (id);
CREATE INDEX IF NOT EXISTS idx_tickets_zone ON tickets (zone);
CREATE INDEX IF NOT EXISTS idx_tickets_status ON tickets (status);
\end{lstlisting}

Esta sintaxis sigue el mecanismo de particionamiento declarativo documentado oficialmente por
Cockroach Labs~\cite{cockroachlabs2024docs}, que exige que la columna de partición sea el primer
componente de la clave primaria para que el optimizador pueda aplicar poda de particiones
(\emph{partition pruning}) en tiempo de planificación de consultas. Esta decisión reemplaza el
diseño original de la Entrega 2, que fragmentaba \texttt{tickets} por
zona geográfica (\texttt{PARTITION BY LIST (zone)}). El cambio fue motivado por un requisito
explícito y no negociable de la guía oficial de la Entrega 3 para el equipo ACC: fragmentación
horizontal de \texttt{tickets} por \texttt{fecha\_apertura}. El razonamiento completo de este
cambio, incluyendo los \emph{trade-offs} aceptados, se documenta en el ADR-0003
(\texttt{docs/adr/0003-sharding-policy.md}) y se resume a continuación.

\subsection{Verificación formal de corrección}

Con $a = \texttt{created\_at}$ y los cuatro puntos de corte del Listado~\ref{lst:init-db}
($c_0 = \text{MINVALUE}$, $c_1 = \text{2026-04-01}$, $c_2 = \text{2026-07-01}$,
$c_3 = \text{2026-10-01}$, $c_4 = \text{MAXVALUE}$), se verifican las tres condiciones del
marco de Özsu y Valduriez~\cite{ozsu2011principles} introducidas en la
Sección~\ref{sec:fundamento-teorico}:

\begin{itemize}[leftmargin=1.5em]
    \item \textbf{Completitud}: la columna \texttt{created\_at} es \texttt{NOT NULL} con valor
        por defecto \texttt{now()}, por lo que todo ticket tiene un valor definido de
        \texttt{created\_at} y, en consecuencia, pertenece a exactamente una de las cuatro
        particiones trimestrales.
    \item \textbf{Reconstrucción}: los rangos
        $[\text{MINVALUE}, c_1) \cup [c_1, c_2) \cup [c_2, c_3) \cup [c_3, \text{MAXVALUE})$
        cubren la totalidad del dominio de \texttt{TIMESTAMPTZ}; la unión \texttt{UNION ALL} de
        las cuatro particiones reconstruye exactamente la tabla \texttt{tickets} sin pérdida de
        filas.
    \item \textbf{Disyunción}: la semántica de \texttt{PARTITION BY RANGE} en CockroachDB define
        cada límite como inclusivo por la izquierda y exclusivo por la derecha
        (\texttt{VALUES FROM ... TO ...}); por construcción, ningún valor de \texttt{created\_at}
        satisface simultáneamente los predicados de dos particiones distintas.
\end{itemize}

Adicionalmente, dado que \texttt{id} deja de ser autosuficiente como punto de acceso —la
aplicación recibe un UUID de ticket sin conocer su \texttt{created\_at} de antemano en los
\emph{endpoints} \texttt{GET}/\texttt{PATCH} sobre \texttt{/api/v1/tickets/\{id\}}—, se añadió el índice único secundario
\texttt{tickets\_id\_key}, usado por \texttt{TicketRepository.findByTicketId(UUID)} para resolver
ese acceso con un único salto índice→tabla, sin necesidad de escanear las cuatro particiones.

\subsection{Colocalización con clientes y limitación arquitectónica}

La guía de la Entrega 3 pide, además de la fragmentación por fecha, colocalizar \texttt{tickets}
con \texttt{clientes} por \texttt{client\_id}, lo que en CockroachDB se implementaría típicamente
con \texttt{INTERLEAVE IN PARENT} o una clave compartida dentro de la misma base de datos. En la
arquitectura de microservicios heredada de la Entrega 2, los clientes son propiedad exclusiva de
\texttt{auth-service} (\texttt{auth\_db.users}), una base de datos físicamente distinta de
\texttt{ticket\_db} — cada microservicio es dueño de su propio esquema, sin acceso cruzado directo
a la base de otro. Forzar una colocalización física en el sentido literal de CockroachDB
requeriría fusionar ambas bases de datos, violando el límite de propiedad de datos entre servicios
que es una decisión arquitectónica deliberada, no un descuido, de la Entrega 2. Se documenta esta
tensión explícitamente en el ADR-0003 en vez de forzar una colocalización artificial: la columna
\texttt{tickets.client\_id} permanece indexada para soportar las consultas
\texttt{findByClientId}/\texttt{findByClientIdAndStatus}, y la resolución de identidad del cliente
ocurre a nivel de aplicación mediante el JWT emitido por \texttt{auth-service}, no a nivel de
almacenamiento físico compartido.

\subsection{Política de replicación}

Como la fragmentación pasó de ser por zona a ser por fecha, ya no existe una correspondencia 1:1
entre una partición y la \emph{locality} de un nodo específico: una partición trimestral contiene
tickets de las tres zonas geográficas mezclados. En consecuencia, la política de replicación
(Listado~\ref{lst:zones}) ya no ancla particiones a nodos por restricción de \emph{locality};
exige uniformemente un factor de replicación 3 sobre toda la tabla, de forma que la pérdida de
cualquier nodo individual del clúster de tres no comprometa la disponibilidad de escritura
(2 de 3 réplicas siguen constituyendo mayoría para Raft~\cite{ongaro2014raft}). Al igual que el
esquema de la Sección anterior, esta política es hoy la migración Flyway
\texttt{V2\_\_configure\_ticket\_zone.sql} (Entregable 5 de la guía de cierre), aplicada
automáticamente después de \texttt{V1} en vez del guion suelto \texttt{db-cluster/config/zones.sql}
que existía hasta entonces —ese archivo se retiró del repositorio, no quedó como una segunda
copia que pudiera desincronizarse de la migración real.

\begin{lstlisting}[language=SQL, caption={Política de replicación uniforme (\texttt{V2\_\_configure\_ticket\_zone.sql})}, label={lst:zones}]
ALTER TABLE tickets CONFIGURE ZONE USING num_replicas = 3;
\end{lstlisting}

\subsection{Trade-offs aceptados}

El rediseño no es gratuito: las consultas filtradas por zona —el filtro más frecuente en la
operación real de un ISP, usado por el RBAC de técnico y por la mayoría de los reportes— dejan de
beneficiarse de poda de particiones y pasan a resolverse mediante \emph{scatter-gather} sobre las
cuatro particiones trimestrales, mientras que antes se resolvían dentro de una sola partición. Este
\emph{trade-off}, junto con el beneficio simétrico obtenido por las consultas de rango temporal
(la mayoría de los reportes reales: ``tickets abiertos hoy'', ``tickets del último trimestre''), se
mide de forma comparativa con \texttt{EXPLAIN ANALYZE} sobre las cinco consultas representativas de
\texttt{db-cluster/scripts/queries\_bench.sql} (lectura por \texttt{id}, rango de fecha, filtro por
zona, agregación de incumplimiento de SLA, y \emph{join} con técnicos), documentadas junto con sus
planes de ejecución en el Anexo correspondiente del repositorio.
